\documentclass[10pt]{article}

\usepackage[margin=0.33in]{geometry}
\usepackage{titlesec}
\usepackage{enumitem}
\usepackage{hyperref}
\usepackage[T1]{fontenc}
\usepackage[scaled]{helvet}
\usepackage[none]{hyphenat}
\sloppy

\hypersetup{hidelinks}
\renewcommand{\familydefault}{\sfdefault}
\setlength{\parindent}{0pt}
\setlength{\parskip}{0pt}
\pagestyle{empty}

% Section formatting (ATS-safe, single-column, no graphics)
\setlength{\parskip}{0pt}
\newcommand{\ressection}[1]{
  \vspace{0.12em}
  {\large\bfseries #1}\par
  \vspace{0.05em}\hrule\vspace{0.12em}
}

% Compact bullet layout for one-page engineering resumes
\setlist[itemize]{
  leftmargin=1.1em,
  itemsep=-1pt,
  topsep=-1pt,
  parsep=0pt,
  partopsep=0pt,
  noitemsep
}

\begin{document}

\begin{center}
{\LARGE {\bfseries Santhosh Sankar}} \\
{\small Aerospace Engineer with expertise in avionics and embedded systems} \\
{\small U.S. Citizen} \\
Ann Arbor, MI \enspace $\vert$ \enspace (734) 846-5255 \enspace $\vert$ \enspace
\href{mailto:ssankar@umich.edu}{ssankar@umich.edu} \enspace $\vert$ \enspace
\href{https://www.linkedin.com/in/santhosh-sankar-3241521b8/}{linkedin.com/in/Tosh-Sankar} \enspace $\vert$ \enspace
\href{https://github.com/ssankar7775}{github.com/ssankar7775}
\end{center}

\ressection{EDUCATION}

{\bfseries University of Michigan}, Ann Arbor, MI \hfill Jan 2026 -- May 2027 (Expected) \\
Master of Engineering in Space Systems \hfill {\bfseries GPA: 4.0/4.0} \\
Climate and Space Sciences and Engineering Department (CLaSP)

\begin{itemize}
  \item Specialization in Controls and Embedded Systems
  \item Relevant Graduate Coursework: Management of Space Systems, Control System Analysis and Design, Astrodynamics
\end{itemize}

{\bfseries University of Michigan}, Ann Arbor, MI \hfill Aug 2021 -- Aug 2024 \\
Bachelor of Science and Engineering in Aerospace Engineering

\ressection{WORK EXPERIENCE}

{\bfseries SpaceX} -- Payload Integration, Launch Engineering \hfill Vandenberg, CA \\
Graduate Engineer \hfill Aug 2025 -- Dec 2025
\begin{itemize}
\item Engineered an additional payload avionics checkout rack, increasing payload processing throughput by 50\% to support increased launch cadence.
\item Developed a simulation environment to validate automated payload avionics checkout sequences, enabling hardware-independent testing and engineer training.
\item Automated payload checkout test plans, reducing required engineering support by 1 hour per launch campaign.
\item Integrated payload avionics hardware supporting Transporter-15, Sentinal 6B and over 20 starlink missions Falcon 9 launch processing operations.
\end{itemize}

\ressection{PROJECT EXPERIENCE}

{\bfseries Michigan Experimental Rocketry Organization (MERO)} \hfill Ann Arbor, MI \\
Avionics Systems Engineer / Responsible Engineer \hfill Jan 2025 -- Present
\begin{itemize}
\item Architected \textbf{MERO} liquid rocket avionics system (\textbf{STM32} flight computer, \textbf{LoRa/SX1276} telemetry, \textbf{24V} power and ground station interfaces), enabling an integrated and scalable \textbf{flight-ready} avionics architecture.
\item Designed and validated \textbf{PCB} power and telemetry subsystems (\textbf{24V} battery bus, regulated \textbf{3.3V} rails, current/voltage monitoring, \textbf{RF} integration) using \textbf{Altium} and circuit simulation tools to ensure reliable system performance.
\item Led avionics \textbf{PDRs} and subsystem coordination, establishing requirements, documentation, and testing workflows to mature \textbf{MERO}'s flight avionics toward \textbf{operational readiness}.
\end{itemize}

{\bfseries CubeSat Flight Lab -- University of Michigan} \hfill Ann Arbor, MI \\
Systems Engineer \hfill Aug 2023 -- Aug 2024
\begin{itemize}
  \item Spearheaded a high-altitude balloon (HAB) flight campaign for a research CubeSat; authored CONOPS, requirements, and integration/test procedures.
  \item Built {\bfseries 6} FlatSat/StratoSat iterations, reducing architecture from 3U 7 kg to 1U 1.2 kg and maturing subsystem integration.
  \item Designed an RP2040-based flight termination unit, reducing subsystem mass by 83\%; validated operation in thermal chamber testing to {\bfseries 110,000 ft altitude}.
  \item Engineered embedded avionics integrating LoRa and GPS; mitigated GPS RF noise via ground-plane integration for robust command, tracking, and mission safety.
\end{itemize}

{\bfseries Michigan Aeronautical Science Association (MASA)} \hfill Ann Arbor, MI \\
Aerodynamics and Recovery Sub-Team Lead \hfill Apr 2022 -- Apr 2023
\begin{itemize}
  \item Led aerodynamic and recovery system development for 23 ft. collegiate liquid bi-propellant (RP-1/LOx) rocket Clementine.
  \item Implemented MBSE workflows for the Lonely Mission demonstrator, introducing Rolleron-based stabilization research to mitigate roll coupling and improve flight verification.
  \item Engineered a transition from pyrotechnic to pneumatic deployment architecture, lowering fracture risk in fiberglass airframes and improving system reliability.
  \item Optimized fin and fin-can geometry using MATLAB (Mastran), FEA, OpenRocket, and CFD; led manufacturing execution including fabrication and composite layups.
\end{itemize}

{\bfseries Human Powered Submarine} \hfill Ann Arbor, MI \\
Propulsion and Hull Design Team Member \hfill Sep 2023 -- Aug 2024
\begin{itemize}
  \item Migrated team CAD workflow to Siemens NX and Teamcenter, enabling large-assembly collaboration and standardizing GD\&T practices across the design team.
  \item Designed a drivetrain tensioner system for gear-and-pulley propulsion, eliminating chain skipping and enabling diver-accessible adjustment during underwater operation.
\end{itemize}

\ressection{TECHNICAL SKILLS}

\textbf{Programming:} Python, C/C++, MATLAB, Simulink \\
\textbf{Embedded Systems:} Microcontrollers (RP2040), Raspberry Pi, embedded firmware, sensor interfaces, telemetry \\
\textbf{Avionics \& Instrumentation:} Oscilloscopes, multimeters, thermal chambers, wind tunnels, spectrum analyzers, Logic Analyzers \\
\textbf{Simulation \& Engineering Tools:} Siemens NX, SolidWorks, Star-CCM+, MATLAB, CFD, FEA, Altium Designer \\
\textbf{Hardware Development:} PCB design, wiring harnesses, composite fabrication, aerospace prototyping \\
\textbf{Project Management:} Requirements, test planning, technical documentation, team leadership
\end{document}